\documentclass[11pt,letterpaper]{article}

\usepackage[spanish,es-noshorthands]{babel}
\usepackage{fontspec}
\setmainfont{Source Serif Pro}
\usepackage{microtype}
\usepackage{geometry}
\usepackage{booktabs}
\usepackage{tabularx}
\usepackage{enumitem}
\usepackage{xcolor}
\usepackage{float}
\usepackage{csquotes}
\usepackage{amsmath}
\usepackage{tikz}
\usetikzlibrary{arrows.meta}
\usepackage{pgfplots}
\pgfplotsset{compat=1.18}
\usepackage[section]{placeins}
\usepackage{xurl}
\usepackage[hidelinks]{hyperref}
\usepackage[backend=biber,style=apa,sorting=nyt]{biblatex}
\DeclareLanguageMapping{spanish}{spanish-apa}

\geometry{margin=2.3cm}
% Los rótulos de las figuras y la bibliografía generan avisos de línea holgada
% que son cosméticos. Se silencian para que `make check` siga siendo sensible a
% los desbordes reales.
\hbadness=10000
\setlength{\emergencystretch}{1em}
\setlength{\parindent}{0pt}
\setlength{\parskip}{0.42em}
\setlist{nosep}
\renewcommand{\topfraction}{0.92}
\renewcommand{\bottomfraction}{0.85}
\renewcommand{\textfraction}{0.06}
\renewcommand{\floatpagefraction}{0.72}
\setcounter{topnumber}{3}
\setcounter{bottomnumber}{2}
\setcounter{totalnumber}{5}

\addbibresource{referencias.bib}
\AtBeginBibliography{\sloppy\small}
\setlength{\bibitemsep}{0.28em}
\setlength{\bibhang}{1.4em}

\newcommand{\pct}{\,\%}
\newcommand{\kusd}{\,k\textdollar}

% Un color por proyecto, más el gris de la gestión del portafolio, el rojo de
% la ruta crítica y el ámbar de las puertas.
\definecolor{cR1}{RGB}{36,102,148}
\definecolor{cR3}{RGB}{56,120,86}
\definecolor{cR4}{RGB}{122,84,150}
\definecolor{cR7}{RGB}{166,92,34}
\definecolor{cDir}{RGB}{132,124,112}
\definecolor{cAviso}{RGB}{176,42,42}
\definecolor{cPuerta}{RGB}{196,140,20}
\definecolor{rBajo}{RGB}{214,232,212}
\definecolor{rMedio}{RGB}{246,224,170}
\definecolor{rAlto}{RGB}{238,178,170}

\title{\vspace{-1.2cm}\textbf{Cuatro proyectos, cuatro puertas}\\
\large Portafolio de I+D de la plataforma de infraestructura como código para
2026-2029, construido a partir de la hoja de ruta tecnológica del primer
programa}
\author{Carlos Luis Rojas Aragonés\\
\small Gestión del Desarrollo Tecnológico: Estrategias y Análisis de Portfolio\\
\small Actividad final 9.5: Portfolio de I+D}
\date{\small 2 de noviembre de 2026}

\begin{document}

\maketitle
\vspace{-0.8em}

\section{El portafolio en una página}
\label{sec:ficha}

En el primer programa construí la hoja de ruta tecnológica de la
infraestructura como código (IaC) de la plataforma en la nube de mi empresa
\autocite{rojas2026hoja}. Esa hoja de ruta tiene ocho proyectos, R1 a R8, para
llevar la plataforma de 2026 a 2029. Para esta actividad elegí cuatro de ellos y
los organicé como un portafolio de I+D, siguiendo los capítulos 16 y 17 de
\textcite{deweck2022}.

\textbf{No agregué proyectos nuevos.} Los proyectos, las figuras de mérito, las
ecuaciones, los precios unitarios, los prototipos y las puertas de decisión
salen de la hoja de ruta. Lo que no estaba allí y tuve que estimar (las
probabilidades de cada puerta y los márgenes de sobrecosto) lo señalo como
estimación propia donde aparece. Todas las cantidades están en miles de dólares
estadounidenses.

\begin{table}[H]
\centering
\footnotesize
\setlength{\tabcolsep}{4pt}
\caption{El portafolio, en corto.}
\label{tab:ficha}
\begin{tabularx}{\textwidth}{@{}>{\raggedright\arraybackslash}p{3.0cm}>{\raggedright\arraybackslash}X@{}}
\toprule
\textbf{Dimensión} & \textbf{Dato} \\
\midrule
Proyectos & R1 Cierre de cobertura (la mitad que falta), R3 Catálogo de módulos
  certificados, R4 Política como código y R7 Interfaz de autoservicio \\
Período & De 2026 T4 a 2029 T2, once trimestres \\
Esfuerzo & 54 meses-persona; 1,6 ingenieros de media y 2,4 en el pico \\
Presupuesto & 473,6 de base más 73,8 de reserva (15,6\pct): \textbf{547,4
  miles de dólares} \\
Valor & VAN esperado de \textbf{727,3} a 2031, con un 5\pct{} de probabilidad
  de quedar por debajo de 456,4 \\
Figuras de mérito & Cobertura de 55 a 86\pct; tiempo de ciclo de 26 a 8 h;
  tasa de fallo de 14 a 6\pct; cumplimiento de 62 a 94\pct;
  reconstrucción de 40 a 11 h \\
Puertas & G1 (2026 T4), G2 (2027 T4), G3 (2028 T3) y G4 (2029 T1), todas
  de la hoja de ruta \\
\bottomrule
\end{tabularx}
\end{table}

Lo que encontré, en tres frases:

\begin{enumerate}
\item \textbf{R1 y R3 son los que pagan.} Solos dan un VAN esperado de 776,8
  a 2031. Toda la hoja de ruta descansa en la cobertura de codificación, y
  estos dos proyectos son los que la mueven.
\item \textbf{R4 y R7 no se pagan solos antes de 2031.} R4 resta 44,7 y R7
  resta 4,8 de valor esperado. Los incluyo igual, porque son los que cumplen
  dos compromisos de la estrategia: el tiempo de ciclo (C3) y el cumplimiento
  de política (C4). El precio de esos dos compromisos es de 49,5 miles de
  dólares de VAN, y prefiero dejarlo escrito antes que esconderlo.
\item \textbf{Las puertas son parte del presupuesto.} Con ellas, la
  probabilidad de pasarse de los 547,4 es del 19,6\pct; si cada proyecto se
  ejecutara completo pase lo que pase, sería del 30\pct.
\end{enumerate}

\section{De la hoja de ruta al portafolio}
\label{sec:seleccion}

\subsection*{El criterio para elegir}

Un portafolio de I+D no es la lista completa de proyectos de una hoja de ruta.
Para \textcite[cap.~16]{deweck2022}, un proyecto de I+D existe para reducir
una incertidumbre técnica y mover una figura de mérito; lo que ya se sabe hacer
es ejecución. Con esa idea usé un criterio simple: \textbf{un proyecto entra en
el portafolio si tiene una pregunta abierta que hay que contestar con un
prototipo antes de comprometer su presupuesto}. La hoja de ruta ya definía esas
preguntas en sus cinco prototipos (P1 a P5), así que el criterio no es
inventado: separa los proyectos que tienen prototipo de los que no.

El Cuadro~\ref{tab:seleccion} aplica el criterio a los ocho proyectos y les
asigna una de las cuatro decisiones con las que el libro revisa una cartera:
detener, cambiar, mantener o iniciar \autocite[fig.~16.17]{deweck2022}.

\begin{table}[H]
\centering
\footnotesize
\setlength{\tabcolsep}{3.5pt}
\caption{Los ocho proyectos de la hoja de ruta y la decisión para cada uno. TRL:
nivel de madurez tecnológica de partida y de llegada \autocite{mankins1995}.
Tipo según \textcite{wheelwright1992}. Costo de la hoja de ruta, sólo
ingeniería.}
\label{tab:seleccion}
\begin{tabularx}{\textwidth}{@{}l>{\raggedright\arraybackslash}p{3.3cm}cl>{\raggedright\arraybackslash}p{1.5cm}r>{\raggedright\arraybackslash}X@{}}
\toprule
\textbf{Cód.} & \textbf{Proyecto} & \textbf{TRL} & \textbf{Prot.} &
  \textbf{Tipo} & \textbf{Costo} & \textbf{Decisión} \\
\midrule
R1 & Cierre de cobertura & 8\,$\to$\,9 & P1, P5 & Derivado & 110,8 &
  \textbf{Mantener}. Va por la mitad; quedan 55,4 \\
R3 & Catálogo de módulos certificados & 7\,$\to$\,9 & P2, P5 & Plataforma &
  142,5 & \textbf{Iniciar} en 2027 T1 \\
R4 & Política como código & 8\,$\to$\,9 & P3 & Derivado & 79,2 &
  \textbf{Iniciar} en 2027 T4 \\
R7 & Interfaz de autoservicio & 6\,$\to$\,9 & P4 & Plataforma & 126,7 &
  \textbf{Iniciar} en 2028 T2, con puerta \\
\midrule
R2 & Custodia del estado & 9\,$\to$\,9 & -- & -- & 47,5 & Cerrado en 2026 \\
R5 & Malla de pruebas & 7\,$\to$\,9 & -- & -- & 95,0 & Cambiar: pasa a
  operación \\
R6 & Sustitución del motor & 8\,$\to$\,9 & -- & -- & 71,2 & Cambiar: pasa a
  operación \\
R8 & Observabilidad & 7\,$\to$\,9 & -- & -- & 63,3 & Cambiar: pasa a
  operación; el portafolio depende de él \\
\bottomrule
\end{tabularx}
\end{table}

R2 ya terminó. R5, R6 y R8 no tienen prototipo porque su incertidumbre es de
costo y de ejecución, no técnica: la misma hoja de ruta decidía \emph{adoptar}
o \emph{comprar} sus herramientas, no construirlas. Los paso al presupuesto de
operación del equipo de plataforma, con sus 229,5 miles de dólares. No
desaparecen, pero no compiten por el dinero de I+D.

\subsection*{Qué tipo de portafolio es}

Ninguno de los cuatro proyectos parte por debajo de TRL 6. Es un portafolio de
desarrollo y maduración, no de investigación básica, y la hoja de ruta explica
por qué: la IaC es una tecnología cuya curva S ya pasó su punto de inflexión, y
la estrategia es cerrar la distancia con lo que ya existe, no apostar por lo que
viene. Por eso las incertidumbres no son de física ni de algoritmos: son de
volumen (cuántos recursos se resisten a importarse), de diseño (si una interfaz
estrecha cubre los casos reales) y de adopción (si los equipos van a usar el
portal). R7 es el único que empieza en TRL 6, y su riesgo es de comportamiento
humano, no técnico.

Los cuatro proyectos no son independientes. R7 necesita el catálogo de R3,
porque el portal sólo expone módulos que existen. R1 y R3 se verifican juntos en
P5, el ensayo de reconstrucción. Y el portafolio depende de R8, que queda
fuera: sin sus mediciones, las puertas deciden con estimaciones y no con datos.

\section{Objetivos y propuesta de valor}
\label{sec:objetivos}

\textbf{Objetivo general:} llevar la plataforma de IaC a los compromisos de la
declaración de estrategia de la hoja de ruta en 2029, con un gasto
comprometido de 547,4 miles de dólares y sin pasar de 2,5 ingenieros a la vez.

El libro propone expresar la propuesta de valor de un proyecto de I+D como
cuánto mueve una figura de mérito y en cuánto tiempo, $\Delta\mathrm{FOM}/\Delta t$
\autocite[ec.~16.1]{deweck2022}. El Cuadro~\ref{tab:fom} lo hace para cada
proyecto, con los valores de 2026 y las metas de 2029 de la hoja de ruta.

\begin{table}[H]
\centering
\footnotesize
\setlength{\tabcolsep}{4pt}
\caption{Propuesta de valor de cada proyecto como $\Delta\mathrm{FOM}/\Delta t$,
con $\Delta t$ de 2026 a 2029. La última columna es el compromiso de la
estrategia (C1 a C5) al que sirve.}
\label{tab:fom}
\begin{tabular}{@{}llrrrl@{}}
\toprule
\textbf{Proyecto} & \textbf{Figura de mérito} & \textbf{2026} &
  \textbf{2029} & $\boldsymbol{\Delta\mathrm{FOM}/\Delta t}$ & \textbf{Comp.} \\
\midrule
R1      & Cobertura de codificación (\%) & 55 & 86 & $+10{,}3$ puntos/año & C1 \\
R1 y R3 & Tiempo de reconstrucción (h)    & 40 & 11 & $-9{,}7$ h/año & C2 \\
R3      & Tasa de fallo del cambio (\%)   & 14 &  7 & $-2{,}3$ puntos/año & -- \\
R4      & Cumplimiento de política (\%)   & 62 & 94 & $+10{,}7$ puntos/año & C4 \\
R7, R4 y R3 & Tiempo de ciclo del cambio (h) & 26 &  8 & $-6{,}0$ h/año & C3 \\
\bottomrule
\end{tabular}
\end{table}

Dicho en una frase por proyecto:

\begin{itemize}
\item \textbf{R1} termina de traer al código los recursos creados a mano. Es la
  figura más importante de la hoja de ruta: explica el 90\pct{} del tiempo de
  reconstrucción de 2026 y es la que más mueve el VAN.
\item \textbf{R3} reemplaza recursos escritos a mano por módulos certificados.
  Por sí solo baja la tasa de fallo del 14 al 7\pct, porque un módulo impide las
  combinaciones de parámetros que nadie validó.
\item \textbf{R4} evalúa las reglas de la organización sobre el plan antes de
  aplicarlo. Sube el cumplimiento automático del 62 al 94\pct{} y quita 4,4
  horas de revisión humana al ciclo.
\item \textbf{R7} pone el catálogo en manos de los equipos de producto. De las
  18 horas que gana el tiempo de ciclo, 12 son suyas, porque es el único que
  elimina la espera en la cola del equipo de plataforma.
\end{itemize}

El compromiso C5, la tasa de deriva, no aparece en el cuadro: depende de la
cobertura de R1 y del intervalo de detección de R8, y R8 está fuera del
portafolio. Es la parte de la estrategia que este portafolio no puede cumplir
solo.

\section{Estructura de división del trabajo}
\label{sec:edt}

El Cuadro~\ref{tab:edt} descompone el portafolio en 19 paquetes de trabajo, cada
uno con un entregable que se puede verificar. El mes-persona se valora a 7,92
miles de dólares, el mismo costo completo de 95 mil dólares anuales por
ingeniero que usa la hoja de ruta. A cada proyecto le sumo el 11,4\pct{} de
otros conceptos (herramientas, servicios y formación), que es la proporción que
la hoja de ruta asignaba sobre la ingeniería.

El elemento 1 es lo único que no estaba en la hoja de ruta. El libro recomienda
no olvidar el costo de gestionar el propio portafolio, y aquí son 3
meses-persona repartidos en todo el período.

\begin{table}[H]
\centering
\footnotesize
\setlength{\tabcolsep}{3.5pt}
\caption{Estructura de división del trabajo. M-p: meses-persona. Costo base en
miles de dólares, con otros conceptos y sin reserva.}
\label{tab:edt}
\begin{tabularx}{\textwidth}{@{}l>{\raggedright\arraybackslash}p{5.4cm}>{\raggedright\arraybackslash}Xrr@{}}
\toprule
\textbf{EDT} & \textbf{Paquete} & \textbf{Entregable} & \textbf{M-p} &
  \textbf{Costo} \\
\midrule
\textbf{1} & \textbf{Gestión del portafolio} & & \textbf{3,0} & \textbf{23,8} \\
1.1 & Gobierno del portafolio y puertas & Actas de G1 a G4 & 1,5 & 11,9 \\
1.2 & Valor ganado y medición de figuras & Informe trimestral de SPI y CPI & 1,5 & 11,9 \\
\addlinespace[2pt]
\textbf{2} & \textbf{R1 Cierre de cobertura (lo que falta)} & & \textbf{7,0} & \textbf{61,7} \\
2.1 & Importación de la cola importable & Recursos importados y conciliados & 3,0 & 26,5 \\
2.2 & Codificación de recursos resistentes & Código para lo que no importa solo & 2,0 & 17,6 \\
2.3 & Conciliación y medición de cobertura & Cobertura medida en 2027 T1 & 1,0 & 8,8 \\
2.4 & P5: reconstrucción cronometrada (parte R1) & Informe del ensayo & 1,0 & 8,8 \\
\addlinespace[2pt]
\textbf{3} & \textbf{R3 Catálogo de módulos certificados} & & \textbf{18,0} & \textbf{158,8} \\
3.1 & P2: módulo de referencia y contrato & Un módulo con tres consumidores & 3,0 & 26,5 \\
3.2 & Módulos de red, cómputo y datos & Módulos publicados en el registro & 7,0 & 61,7 \\
3.3 & Certificación: firma, versiones y registro & Contrato de interfaz firmado & 4,0 & 35,3 \\
3.4 & Migración de consumidores & Penetración $\kappa$ medida & 3,0 & 26,5 \\
3.5 & P5: reconstrucción cronometrada (parte R3) & Informe del ensayo & 1,0 & 8,8 \\
\addlinespace[2pt]
\textbf{4} & \textbf{R4 Política como código} & & \textbf{10,0} & \textbf{88,2} \\
4.1 & P3: evaluación en modo aviso & Tasa de falsos positivos medida & 2,0 & 17,6 \\
4.2 & Biblioteca de reglas propias & Reglas versionadas & 4,0 & 35,3 \\
4.3 & Modo bloqueo y régimen de excepciones & Política bloqueante en producción & 3,0 & 26,5 \\
4.4 & Medición de falsos positivos & Cumplimiento $\gamma$ medido & 1,0 & 8,8 \\
\addlinespace[2pt]
\textbf{5} & \textbf{R7 Interfaz de autoservicio} & & \textbf{16,0} & \textbf{141,1} \\
5.1 & P4: portal, tres plantillas, diez equipos & Adopción $\sigma$ medida & 3,0 & 26,5 \\
5.2 & Plantillas sobre el catálogo & Plantillas publicadas & 6,0 & 52,9 \\
5.3 & Integración con petición y aprobación & Flujo de aprobación en producción & 4,0 & 35,3 \\
5.4 & Despliegue a todos los equipos & Portal abierto a toda la empresa & 3,0 & 26,5 \\
\midrule
 & \textbf{Total} & & \textbf{54,0} & \textbf{473,6} \\
\bottomrule
\end{tabularx}
\end{table}

\section{Calendario}
\label{sec:calendario}

La Figura~\ref{fig:gantt} programa los paquetes en los once trimestres del
portafolio. Respeta las fechas de la hoja de ruta: R3 en 2027, R4 en 2028 y R7
entre 2028 y 2029, con los prototipos adelantados a su proyecto para que cada
puerta decida con una medición y no con una suposición.

\begin{figure}[H]
\centering
\begin{tikzpicture}[x=0.97cm, y=0.43cm, font=\scriptsize]
  \foreach \i in {0,...,11} \draw[gray!20] (\i,0.9) -- (\i,-22.0);
  \foreach \i in {1,5,9} \draw[gray!55, thick] (\i,0.9) -- (\i,-22.0);
  \foreach \i/\a in {0.5/2026, 3/2027, 7/2028, 10/2029}
    \node[anchor=south, font=\scriptsize\bfseries] at (\i,1.0) {\a};
  \foreach \i/\t in {0.5/T4, 1.5/T1, 2.5/T2, 3.5/T3, 4.5/T4, 5.5/T1, 6.5/T2,
                     7.5/T3, 8.5/T4, 9.5/T1, 10.5/T2}
    \node[anchor=north, text=gray!50!black, font=\tiny] at (\i,-22.1) {\t};
  % fila / inicio / fin / color / rótulo / ruta crítica
  \foreach \f/\a/\b/\c/\r/\k in {
     0/0/11/cDir/{1 Gestión del portafolio}/0,
     2/0/1/cR1/{2.1 Importación de la cola importable}/0,
     3/0/2/cR1/{2.2 Codificación de recursos resistentes}/0,
     4/1/2/cR1/{2.3 Conciliación y medición}/0,
     5/8/9/cR1/{2.4 P5: reconstrucción (R1)}/0,
     7/1/2/cR3/{3.1 P2: módulo de referencia}/1,
     8/2/4/cR3/{3.2 Módulos de red, cómputo y datos}/1,
     9/3/5/cR3/{3.3 Certificación}/1,
     10/4/5/cR3/{3.4 Migración de consumidores}/0,
     11/8/9/cR3/{3.5 P5: reconstrucción (R3)}/0,
     13/4/5/cR4/{4.1 P3: política en modo aviso}/0,
     14/5/6/cR4/{4.2 Biblioteca de reglas propias}/0,
     15/5/7/cR4/{4.3 Modo bloqueo y excepciones}/0,
     16/6/7/cR4/{4.4 Medición de falsos positivos}/0,
     18/6/8/cR7/{5.1 P4: portal y medición de $\sigma$}/1,
     19/8/10/cR7/{5.2 Plantillas sobre el catálogo}/1,
     20/9/10/cR7/{5.3 Integración con petición}/1,
     21/10/11/cR7/{5.4 Despliegue a todos los equipos}/1} {
    \fill[\c!26, draw=\c, thick, rounded corners=1.2pt]
      (\a+0.04,-\f+0.32) rectangle (\b-0.04,-\f-0.32);
    \ifnum\k=1
      \draw[cAviso, very thick, rounded corners=1.2pt]
        (\a+0.04,-\f+0.32) rectangle (\b-0.04,-\f-0.32);
    \fi
    \node[anchor=east, font=\tiny, text=gray!15!black] at (-0.12,-\f) {\r};
  }
  % holgura de la ruta crítica en 2028 T1
  \draw[cAviso, thick, densely dotted, -{Stealth[length=4pt]}]
    (5.0,-9) -- (5.0,-17.4) -- (5.96,-17.4) -- (5.96,-17.65);
  \node[cAviso, font=\tiny, anchor=east, align=right] at (4.94,-12.0)
    {holgura de\\un trimestre};
  % puertas
  \foreach \g/\x in {G1/1, G2/5, G3/8, G4/10} {
    \draw[cPuerta!80!black, thick, densely dashed] (\x,0.85) -- (\x,-22.0);
    \node[fill=white, draw=cPuerta!80!black, thick, rounded corners=1.4pt,
          inner sep=1.4pt, font=\tiny\bfseries] at (\x,0.4) {\g};
  }
\end{tikzpicture}
\caption{Calendario del portafolio, de 2026 T4 a 2029 T2. Un color por
proyecto y gris para la gestión. Contorno rojo: ruta crítica. Líneas ámbar:
puertas de decisión de la hoja de ruta. Elaboración propia.}
\label{fig:gantt}
\end{figure}

La ruta crítica es \textbf{3.1, 3.2, 3.3, 5.1, G3, 5.2, 5.3 y 5.4}: el
catálogo, el prototipo del portal, la puerta de adopción y el autoservicio
completo. Fija el final del portafolio en 2029 T2 y tiene una sola holgura, el
primer trimestre de 2028, entre la certificación del catálogo y el prototipo
P4. Si la certificación se atrasa un trimestre, el portafolio termina a tiempo;
si se atrasa dos, R7 se corre y con él el compromiso del tiempo de ciclo.

El Cuadro~\ref{tab:puertas} resume las cuatro puertas. Son las de la hoja de
ruta, con su criterio numérico y la consecuencia decidida de antemano, como en
el modelo Stage-Gate de \textcite{cooper1990}. G1 se decide este mismo
trimestre.

\begin{table}[H]
\centering
\footnotesize
\setlength{\tabcolsep}{4pt}
\caption{Puertas de decisión del portafolio.}
\label{tab:puertas}
\begin{tabularx}{\textwidth}{@{}lll>{\raggedright\arraybackslash}X@{}}
\toprule
\textbf{Puerta} & \textbf{Fecha} & \textbf{Evidencia} &
  \textbf{Criterio y consecuencia de no superarla} \\
\midrule
G1 & 2026 T4 & P1 & Al menos el 70\pct{} de los recursos se importa sin
  intervención. Si no: la meta de cobertura baja y se recalcula el valor del
  portafolio antes de iniciar R3. \\
G2 & 2027 T4 & P2, P3 & El módulo de referencia sirve a tres consumidores sin
  bifurcarse, y la política tiene menos del 15\pct{} de falsos positivos. Si
  no: R4 se queda en modo aviso, no se construye el modo bloqueo (4.3) y el
  cumplimiento se queda cerca del 70\pct. \\
G3 & 2028 T3 & P4 & Adopción $\sigma \geq 0{,}60$ con diez equipos reales.
  Si no: \textbf{R7 se cancela} habiendo gastado 26,5 de sus 141,1. \\
G4 & 2029 T1 & P5 & Reconstrucción completa en menos de 16 horas. Si no: no
  se declara cumplido el compromiso C2 y se abre un proyecto de continuidad
  fuera del portafolio. \\
\bottomrule
\end{tabularx}
\end{table}

\section{Presupuesto}
\label{sec:presupuesto}

\subsection*{Por proyecto y por trimestre}

El Cuadro~\ref{tab:presupuesto} reparte el costo de la EDT en el calendario, y
la Figura~\ref{fig:arena} lo muestra como gráfico de arena y como curva
acumulada, que son las dos vistas que propone el libro
\autocite[sec.~16.2.3]{deweck2022}.

\begin{table}[H]
\centering
\footnotesize
\setlength{\tabcolsep}{2.8pt}
\caption{Presupuesto base por proyecto y trimestre, en miles de dólares, sin
reserva. Las sumas pueden diferir en una décima por redondeo.}
\label{tab:presupuesto}
\begin{tabular}{@{}l*{11}{r}r@{}}
\toprule
 & \multicolumn{1}{c}{\textbf{2026}} & \multicolumn{4}{c}{\textbf{2027}} &
   \multicolumn{4}{c}{\textbf{2028}} & \multicolumn{2}{c}{\textbf{2029}} & \\
\cmidrule(lr){2-2}\cmidrule(lr){3-6}\cmidrule(lr){7-10}\cmidrule(lr){11-12}
\textbf{EDT} & T4 & T1 & T2 & T3 & T4 & T1 & T2 & T3 & T4 & T1 & T2 &
  \textbf{Total} \\
\midrule
1 Gestión & 2,2 & 2,2 & 2,2 & 2,2 & 2,2 & 2,2 & 2,2 & 2,2 & 2,2 & 2,2 & 2,2 & 23,8 \\
2 R1      & 35,3 & 17,6 & & & & & & & 8,8 & & & 61,7 \\
3 R3      & & 26,5 & 30,9 & 48,5 & 44,1 & & & & 8,8 & & & 158,8 \\
4 R4      & & & & & 17,6 & 48,5 & 22,0 & & & & & 88,2 \\
5 R7      & & & & & & & 13,2 & 13,2 & 26,5 & 61,7 & 26,5 & 141,1 \\
\midrule
\textbf{Trimestre} & 37,4 & 46,3 & 33,0 & 50,7 & \textbf{63,9} & 50,7 & 37,4 &
  15,4 & 46,3 & \textbf{63,9} & 28,6 & \textbf{473,6} \\
\textbf{Acumulado} & 37,4 & 83,7 & 116,7 & 167,4 & 231,3 & 282,0 & 319,4 &
  334,8 & 381,0 & 444,9 & 473,6 & \\
\bottomrule
\end{tabular}
\end{table}

\begin{figure}[H]
\centering
\begin{tikzpicture}
\begin{axis}[
  name=arena, width=8.4cm, height=6.0cm,
  ybar stacked, bar width=10pt,
  xmin=-0.6, xmax=10.6, ymin=0, ymax=72,
  xtick={0,...,10},
  xticklabels={26T4,27T1,T2,T3,T4,28T1,T2,T3,T4,29T1,T2},
  xticklabel style={font=\tiny, rotate=60, anchor=east},
  yticklabel style={font=\scriptsize},
  ylabel={\footnotesize Miles de dólares por trimestre},
  ymajorgrids, grid style={gray!18}, axis lines=left,
  legend style={font=\tiny, at={(0.5,-0.30)}, anchor=north,
                draw=gray!40, fill=white, legend columns=5},
  area legend,
  title={\footnotesize Gráfico de arena},
]
  \addplot[fill=cDir!45, draw=cDir] coordinates {(0,2.2)(1,2.2)(2,2.2)(3,2.2)(4,2.2)(5,2.2)(6,2.2)(7,2.2)(8,2.2)(9,2.2)(10,2.2)};
  \addlegendentry{Gestión}
  \addplot[fill=cR1!35, draw=cR1] coordinates {(0,35.3)(1,17.6)(2,0)(3,0)(4,0)(5,0)(6,0)(7,0)(8,8.8)(9,0)(10,0)};
  \addlegendentry{R1}
  \addplot[fill=cR3!35, draw=cR3] coordinates {(0,0)(1,26.5)(2,30.9)(3,48.5)(4,44.1)(5,0)(6,0)(7,0)(8,8.8)(9,0)(10,0)};
  \addlegendentry{R3}
  \addplot[fill=cR4!35, draw=cR4] coordinates {(0,0)(1,0)(2,0)(3,0)(4,17.6)(5,48.5)(6,22.0)(7,0)(8,0)(9,0)(10,0)};
  \addlegendentry{R4}
  \addplot[fill=cR7!35, draw=cR7] coordinates {(0,0)(1,0)(2,0)(3,0)(4,0)(5,0)(6,13.2)(7,13.2)(8,26.5)(9,61.7)(10,26.5)};
  \addlegendentry{R7}
\end{axis}
\begin{axis}[
  at={(arena.east)}, anchor=west, xshift=1.1cm,
  width=7.4cm, height=6.0cm,
  xmin=0, xmax=11, ymin=0, ymax=680,
  xtick={0,1,5,9,11},
  xticklabels={26T4,27T1,28T1,29T1,fin},
  xticklabel style={font=\tiny},
  yticklabel style={font=\scriptsize},
  ytick={0,100,200,300,400,500,600},
  ymajorgrids, grid style={gray!18}, axis lines=left,
  title={\footnotesize Curva acumulada},
  clip=false,
]
  \addplot[cR1, very thick, mark=*, mark size=1.4pt] coordinates {
    (0,0)(1,37.4)(2,83.7)(3,116.7)(4,167.4)(5,231.3)(6,282.0)(7,319.4)
    (8,334.8)(9,381.0)(10,444.9)(11,473.6)};
  \draw[cAviso, thick] (axis cs:0,547.4) -- (axis cs:11,547.4)
    node[pos=0.03, anchor=south west, font=\tiny] {Base más reserva: 547,4};
  \draw[gray!60!black, thick, densely dashed] (axis cs:0,630.9) -- (axis cs:11,630.9)
    node[pos=0.03, anchor=south west, font=\tiny] {P80 con el sobrecosto del sector: 630,9};
  \node[font=\tiny, text=cR1, anchor=west, align=left] at (axis cs:6.4,200)
    {2028 T3: el gasto\\casi se detiene\\mientras se mide $\sigma$};
\end{axis}
\end{tikzpicture}
\caption{Gasto del portafolio por trimestre y acumulado. Los dos picos, 2027 T4
y 2029 T1, coinciden con la mayor concurrencia de paquetes. La pausa de 2028 T3
es la medición de la adopción antes de G3. Elaboración propia.}
\label{fig:arena}
\end{figure}

La curva acumulada no tiene la forma de S que el libro describe como típica:
tiene un escalón en 2028 T3. Es deliberado. Mientras los diez equipos usan el
prototipo del portal, el gasto casi se detiene, y sólo se compromete el resto de
R7 (114,6) si la adopción supera el umbral.

\subsection*{Reserva y riesgo de costo}

El libro recomienda una reserva de gestión de entre el 10 y el 30\pct{}
\autocite[sec.~16.2.3]{deweck2022}. La gradué por madurez, porque el riesgo no es
igual en todos los proyectos: 10\pct{} para la gestión y para R1, que ya va por
la mitad; 15\pct{} para R3 y R4; y 20\pct{} para R7, el de menor TRL.

\begin{table}[H]
\centering
\footnotesize
\setlength{\tabcolsep}{6pt}
\caption{Reserva de gestión por proyecto, en miles de dólares.}
\label{tab:reserva}
\begin{tabular}{@{}lrrrr@{}}
\toprule
\textbf{Proyecto} & \textbf{Base} & \textbf{Tasa} & \textbf{Reserva} &
  \textbf{Total} \\
\midrule
1 Gestión & 23,8 & 10\pct & 2,4 & 26,1 \\
R1 & 61,7 & 10\pct & 6,2 & 67,9 \\
R3 & 158,8 & 15\pct & 23,8 & 182,6 \\
R4 & 88,2 & 15\pct & 13,2 & 101,4 \\
R7 & 141,1 & 20\pct & 28,2 & 169,3 \\
\midrule
\textbf{Portafolio} & \textbf{473,6} & \textbf{15,6\pct} & \textbf{73,8} &
  \textbf{547,4} \\
\bottomrule
\end{tabular}
\end{table}

Para saber si alcanza, simulé el costo real de cada proyecto 40.000 veces con
un sobrecosto aleatorio (estimación propia: mediana del 6\pct{} en R1, 8\pct{} en
R3, 10\pct{} en R4 y 12\pct{} en R7, con cola larga) y apliqué las puertas: un
proyecto detenido sólo cuesta lo gastado hasta su puerta. El resultado es un
costo esperado de 483,6, un P80 de 546,5 y una probabilidad del 19,6\pct{} de
pasarse de los 547,4. El libro pide planificar para el P50 o el P80
\autocite[sec.~16.3]{deweck2022}, y el presupuesto queda justo en el P80.

Sin puertas, es decir, si cada proyecto se ejecutara completo pase lo que pase,
el P80 subiría a 564,1 y la probabilidad de pasarse al 30\pct. \textbf{La
reserva sólo alcanza porque existen las puertas.} Saltarse una es gastar
reserva que no existe.

\subsection*{Capacidad y seguimiento}

Los 54 meses-persona en 33 meses dan 1,6 ingenieros de media, con picos de 2,4
en 2027 T4 y 2029 T1. La hoja de ruta estimaba un equipo de 2,6 ingenieros, así
que el portafolio cabe, pero con poco margen en los picos.

El seguimiento será por valor ganado: la fila acumulada del
Cuadro~\ref{tab:presupuesto} es la línea base, y cada trimestre se reportan el
índice de plazo, $\mathrm{SPI} = \mathrm{BCWP}/\mathrm{BCWS}$, y el de costo,
$\mathrm{CPI} = \mathrm{BCWP}/\mathrm{ACWP}$ \autocite[sec.~16.3]{deweck2022}. Para
evitar el error más común, declarar tareas al 90\pct{} que van por la mitad,
sólo acredito el trabajo al 0, 50 o 100\pct, y el 100\pct{} únicamente cuando
alguien de fuera del proyecto verifica el entregable del Cuadro~\ref{tab:edt}.

\section{Matriz de riesgo}
\label{sec:riesgo}

Uso la matriz de cinco por cinco de la NASA que reproduce el libro
\autocite[fig.~16.9]{deweck2022}, con criterios de impacto adaptados al tamaño
de este portafolio (Cuadro~\ref{tab:criterios}). La Figura~\ref{fig:matriz}
ubica los ocho riesgos y el Cuadro~\ref{tab:riesgos} dice qué hago con cada
uno.

\begin{table}[H]
\centering
\footnotesize
\setlength{\tabcolsep}{3.5pt}
\caption{Criterios de probabilidad e impacto. El impacto de un riesgo es el
mayor que alcanza en cualquiera de las tres columnas.}
\label{tab:criterios}
\begin{tabularx}{\textwidth}{@{}c>{\raggedright\arraybackslash}p{2.1cm}>{\raggedright\arraybackslash}X>{\raggedright\arraybackslash}p{2.3cm}>{\raggedright\arraybackslash}p{3.6cm}@{}}
\toprule
\textbf{Nivel} & \textbf{Probabilidad} & \textbf{Técnico} & \textbf{Costo o
  VAN} & \textbf{Plazo} \\
\midrule
5 & Más del 75\pct  & Se pierde C1, el compromiso principal & Más de 100\kusd &
  El portafolio pasa de 2029 \\
4 & 50 a 75\pct     & Se pierde otro compromiso o se cancela un proyecto & 50 a
  100\kusd & Más de dos trimestres en la ruta crítica \\
3 & 30 a 50\pct     & Una figura de mérito queda bajo su meta & 20 a 50\kusd &
  Se consume la holgura de 2028 T1 \\
2 & 10 a 30\pct     & Se reduce el margen de una figura & 5 a 20\kusd & Un
  trimestre fuera de la ruta crítica \\
1 & Menos del 10\pct & Sin efecto apreciable & Menos de 5\kusd & Sin efecto \\
\bottomrule
\end{tabularx}
\end{table}

\begin{figure}[H]
\centering
\begin{tikzpicture}[x=1.55cm, y=0.95cm]
  \foreach \p/\i/\c in {
    1/1/rBajo, 1/2/rBajo, 1/3/rBajo, 1/4/rMedio, 1/5/rMedio,
    2/1/rBajo, 2/2/rBajo, 2/3/rMedio, 2/4/rMedio, 2/5/rAlto,
    3/1/rBajo, 3/2/rMedio, 3/3/rMedio, 3/4/rAlto, 3/5/rAlto,
    4/1/rBajo, 4/2/rMedio, 4/3/rAlto, 4/4/rAlto, 4/5/rAlto,
    5/1/rMedio, 5/2/rMedio, 5/3/rAlto, 5/4/rAlto, 5/5/rAlto} {
    \fill[\c, draw=white, line width=1.2pt] (\i-1,\p-1) rectangle (\i,\p);
  }
  \foreach \i/\t in {1/Despreciable, 2/Marginal, 3/Moderado, 4/Crítico, 5/Catastrófico}
    \node[font=\tiny, anchor=north] at (\i-0.5,-0.05) {\i\ \t};
  \foreach \p/\t in {1/{Improbable}, 2/{Poco probable}, 3/{Probable}, 4/{Muy probable}, 5/{Casi seguro}}
    \node[font=\tiny, anchor=east] at (-0.06,\p-0.5) {\t\ \p};
  \node[font=\scriptsize\bfseries] at (2.5,-0.62) {Impacto};
  \node[font=\scriptsize\bfseries, rotate=90] at (-1.35,2.5) {Probabilidad};
  \tikzset{r/.style={circle, fill=white, draw=gray!30!black, inner sep=0.9pt,
                     font=\tiny\bfseries, minimum size=13pt}}
  \node[r] at (4.5,1.5) {K1};
  \node[r] at (3.3,1.5) {K2};  \node[r] at (3.7,1.5) {K3};
  \node[r] at (2.2,1.5) {K4};  \node[r] at (2.5,1.5) {K5};  \node[r] at (2.8,1.5) {K6};
  \node[r] at (2.3,2.5) {K7};  \node[r] at (2.7,2.5) {K8};
\end{tikzpicture}
\caption{Matriz de riesgo del portafolio. Verde: aceptar y vigilar. Ámbar:
mitigar con responsable y disparador. Rojo: mitigar antes de iniciar el paquete
afectado. Adaptada de \textcite[fig.~16.9]{deweck2022}.}
\label{fig:matriz}
\end{figure}

\begin{table}[H]
\centering
\footnotesize
\setlength{\tabcolsep}{3pt}
\caption{Registro de riesgos. P: probabilidad. I: impacto. Las probabilidades
de K1, K2, K3 y K4 son las del modelo de la sección~\ref{sec:valor}.}
\label{tab:riesgos}
\begin{tabularx}{\textwidth}{@{}l>{\raggedright\arraybackslash}p{4.6cm}cc>{\raggedright\arraybackslash}X@{}}
\toprule
\textbf{Id.} & \textbf{Riesgo} & \textbf{P} & \textbf{I} &
  \textbf{Respuesta y disparador} \\
\midrule
K1 & La cobertura se estanca en el 74\pct{} (25\pct) & 2 & 5 & Mitigar: G1 lo
  detecta este trimestre y 2.2 ataca los recursos que no importan solos.
  Disparador: menos del 70\pct{} importable en P1. \\
K2 & La adopción $\sigma$ no llega a 0,60 en G3 (26,5\pct) & 2 & 4 & Aceptar
  con puerta: P4 la mide antes de gastar el 81\pct{} de R7, y la cancelación
  está presupuestada. \\
K3 & Falsos positivos de política por encima del 15\pct{} (20\pct) & 2 & 4 &
  Aceptar con puerta: en G2, R4 se queda en modo aviso y se ahorra 4.3. \\
K4 & El sobrecosto supera la reserva (19,6\pct) & 2 & 3 & Mitigar: reserva
  graduada y puertas antes de la mitad del gasto de R4 y R7. \\
K5 & El catálogo se llena de variantes y la interfaz deja de ser estrecha
  (18\pct) & 2 & 3 & Mitigar: contrato de interfaz en P2 y criterio de G2 de
  tres consumidores sin bifurcación. \\
K6 & La reconstrucción de P5 pasa de 16 horas & 2 & 3 & Aceptar con puerta:
  G4 abre un proyecto de continuidad fuera del portafolio. \\
K7 & Los picos de 2,4 ingenieros chocan con trabajo de clientes o de
  operación & 3 & 3 & Mitigar: la holgura de 2028 T1 absorbe un trimestre; 3.4
  es el primer paquete que se mueve. \\
K8 & R8, fuera del portafolio, no entrega a tiempo las mediciones & 3 & 3 &
  Mitigar: R8 es dependencia declarada; sin sus datos, una puerta se aplaza,
  no se aprueba. \\
\bottomrule
\end{tabularx}
\end{table}

El único riesgo rojo es K1, y no es casualidad: la cobertura es la variable de
la que depende casi todo el valor (sección~\ref{sec:valor}). Por eso G1 es la
primera puerta y se decide antes de comprometer un solo dólar de R3. Los riesgos
K2 y K3 son críticos pero quedan en ámbar porque tienen puerta: si ocurren,
cuestan poco, porque se descubren antes de gastar.

\section{Valor y riesgo del portafolio}
\label{sec:valor}

\subsection*{El modelo}

Para valorar el portafolio usé el modelo financiero de la hoja de ruta sin
cambiar sus ecuaciones. Cada beneficio sale de una figura de mérito, y una
mejora que no se puede escribir como función de una figura no entra:

\begin{itemize}
\item \textbf{B1, operación manual evitada:} $N(t)\,[C(t) - 0{,}55]\times 0{,}4$,
  con $N$ el número de recursos y $C$ la cobertura. Lo mueve R1.
\item \textbf{B2, retrabajo evitado:} $n_c(t)\,[0{,}14 - f(t)]\times 1{,}6$,
  con $f = f_0(1-\pi\gamma)(1-\theta\kappa)$, $\pi = 0{,}354$ y
  $\theta = 0{,}688$. Lo mueven R3 con $\kappa$ y R4 con $\gamma$.
\item \textbf{B3, capacidad de plataforma liberada:} 2,1 horas de acompañamiento
  por cambio, a 60 dólares la hora, que suprimen R3 y R7.
\end{itemize}

La cuarta corriente de la hoja de ruta, el menor costo unitario, viene de
sustituir el motor (R6) y queda fuera. Las incertidumbres son los precios
unitarios ($\pm 30$\pct), el crecimiento del parque (5 a 22\pct{} anual), el
resultado de las tres puertas que pueden fallar (G1, G2 y G3, con las
probabilidades del Cuadro~\ref{tab:riesgos}) y el sobrecosto de cada proyecto.
La adopción $\sigma$ sigue una distribución beta de media 0,68, que da un
26,5\pct{} de probabilidad de no pasar G3. Descuento al 9\pct{} hasta 2031,
igual que la hoja de ruta, y no cuento el costo ya gastado de R1, que es costo
hundido. El cálculo está en \texttt{modelo/portafolio.py} y se reproduce
exactamente con la semilla fija.

\subsection*{Seis portafolios posibles}

R1 se mantiene en todos los casos, porque ya va por la mitad y la estrategia
dice que, si sólo se puede cumplir un compromiso, es el de cobertura. Sobre esa
base, las combinaciones posibles de R3, R4 y R7 son seis (R7 no puede ir sin
R3). El Cuadro~\ref{tab:carteras} y la Figura~\ref{fig:frontera} las comparan
por valor esperado y por riesgo, medido como desviación estándar del VAN, que
es la lógica de media y varianza de \textcite{markowitz1952} que el libro aplica
a las carteras de I+D \autocite[sec.~16.5]{deweck2022}.

\begin{table}[H]
\centering
\footnotesize
\setlength{\tabcolsep}{5pt}
\caption{Valor y riesgo de los seis portafolios posibles, en miles de dólares,
a 2031. Ninguno tiene probabilidad apreciable de VAN negativo.}
\label{tab:carteras}
\begin{tabular}{@{}lrrrl@{}}
\toprule
\textbf{Portafolio} & \textbf{Costo total} & \textbf{VAN esperado} &
  \textbf{Desv. est.} & \textbf{Posición} \\
\midrule
R1            &  94,0 & 552,1 & 148,5 & Frontera \\
R1 + R3       & 276,6 & \textbf{776,8} & 158,8 & Frontera, máximo VAN \\
R1 + R4       & 195,5 & 535,3 & 150,0 & Dominado \\
R1 + R3 + R4  & 378,0 & 732,1 & 159,8 & Dominado \\
R1 + R3 + R7  & 446,0 & 772,0 & 163,3 & Dominado \\
\textbf{R1 + R3 + R4 + R7} & \textbf{547,4} & \textbf{727,3} & \textbf{164,4} &
  Dominado; \textbf{recomendado} \\
\bottomrule
\end{tabular}
\end{table}

\begin{figure}[H]
\centering
\begin{tikzpicture}
\begin{axis}[
  width=11.5cm, height=6.4cm,
  xmin=146, xmax=167, ymin=500, ymax=810,
  xlabel={\footnotesize Riesgo: desviación estándar del VAN, miles de dólares},
  ylabel={\footnotesize VAN esperado a 2031},
  xticklabel style={font=\scriptsize}, yticklabel style={font=\scriptsize},
  grid=major, grid style={gray!18}, axis lines=left, clip=false,
]
  \addplot[cR1, thick, densely dashed] coordinates {(148.5,552.1) (158.8,776.8)};
  \addplot[only marks, mark=*, mark size=2.4pt, cR1] coordinates {(148.5,552.1) (158.8,776.8)};
  \addplot[only marks, mark=o, mark size=2.4pt, gray!60!black, thick] coordinates {
    (150.0,535.3) (159.8,732.1) (163.3,772.0)};
  \addplot[only marks, mark=square*, mark size=2.8pt, cR7] coordinates {(164.4,727.3)};
  \node[font=\tiny, anchor=west] at (axis cs:148.9,552.1) {R1};
  \node[font=\tiny, anchor=south east] at (axis cs:158.6,778) {R1+R3};
  \node[font=\tiny, anchor=west] at (axis cs:150.4,535.3) {R1+R4};
  \node[font=\tiny, anchor=north] at (axis cs:159.8,728) {R1+R3+R4};
  \node[font=\tiny, anchor=south] at (axis cs:163.3,775) {R1+R3+R7};
  \node[font=\tiny\bfseries, text=cR7, anchor=north] at (axis cs:164.4,722)
    {Los cuatro};
  \draw[cR7, thin, -{Stealth[length=4pt]}] (axis cs:159.6,770) --
    (axis cs:164.0,731);
  \node[font=\tiny, text=cR7, anchor=west, align=left] at (axis cs:160.6,620)
    {Precio de C3 y C4:\\49,5 de VAN};
\end{axis}
\end{tikzpicture}
\caption{Valor esperado contra riesgo de los seis portafolios. Puntos llenos:
frontera eficiente. Círculos: portafolios dominados. Cuadrado: el portafolio
recomendado. Elaboración propia.}
\label{fig:frontera}
\end{figure}

\subsection*{Qué aporta cada proyecto}

Quitando un proyecto a la vez del portafolio completo se obtiene lo que aporta
cada uno:

\begin{itemize}
\item \textbf{R1:} $+573{,}4$. Es el proyecto que sostiene todo el portafolio.
\item \textbf{R3, junto con R7, que depende de él:} $+192{,}0$, con una
  probabilidad de pérdida del 0,1\pct.
\item \textbf{R7:} $-4{,}8$ a 2031, con un 60\pct{} de probabilidad de restar.
  Si el horizonte se alarga a 2033, aporta $+73{,}5$ y la probabilidad de
  restar baja al 27\pct. Es un proyecto que llega tarde al modelo, no un mal
  proyecto.
\item \textbf{R4:} $-44{,}7$ a 2031 y $-25{,}0$ a 2033. Resta en casi todos los
  escenarios, porque su beneficio en dinero es sólo la reducción de la tasa de
  fallo, y el catálogo ya hace la mayor parte de ese trabajo.
\end{itemize}

\subsection*{La decisión}

Si sólo mirara el dinero, el portafolio sería R1 y R3. Recomiendo los cuatro,
sabiendo que es un portafolio dominado, por dos razones que el modelo no
captura:

\begin{enumerate}
\item \textbf{R7 es el tiempo de ciclo.} Sin él, el tiempo de ciclo de 2029 se
  queda en 20 horas en lugar de 8, según la ecuación de la hoja de ruta, y el
  compromiso C3 no se cumple. El tiempo de ciclo es lo que el negocio percibe
  como plazo de salida al mercado, y el modelo no le pone precio. Además, G3
  limita la pérdida: si la adopción no llega, R7 se cancela habiendo gastado
  26,5.
\item \textbf{R4 es el cumplimiento demostrable.} La motivación E3 de la hoja de
  ruta es que cada auditoría deje de ser un proyecto. El modelo no cuenta el
  costo de una auditoría ni el de una excepción manual, así que R4 aparece como
  pérdida. Si la empresa no tiene una obligación real de cumplimiento, R4 es el
  primer proyecto que yo quitaría, y el ahorro serían 101,4 miles de dólares
  del presupuesto.
\end{enumerate}

Los 49,5 miles de dólares de VAN que se dejan sobre la mesa son el precio de
esos dos compromisos, y es una decisión que tiene que tomar la dirección con el
número a la vista, no el equipo técnico por su cuenta.

\section{Contraste con el sector}
\label{sec:contraste}

La consigna pide contrastar el proyecto con profesionales del sector. Lo hice
con lo que el sector publica: informes de adopción y datos de sobrecostos.
\textbf{No hice entrevistas}; al final de la sección dejo las preguntas que le
haría a un colega. El Cuadro~\ref{tab:contraste} compara cada supuesto con la
evidencia externa.

\begin{table}[H]
\centering
\footnotesize
\setlength{\tabcolsep}{3.5pt}
\caption{Supuestos del portafolio frente a la evidencia del sector. (P) indica
dato de un proveedor con interés comercial.}
\label{tab:contraste}
\begin{tabularx}{\textwidth}{@{}>{\raggedright\arraybackslash}p{3.1cm}>{\raggedright\arraybackslash}X>{\raggedright\arraybackslash}p{4.6cm}@{}}
\toprule
\textbf{Supuesto} & \textbf{Evidencia} & \textbf{Veredicto} \\
\midrule
Cobertura del 55 al 86\pct{} en tres años
& El 89\pct{} de las organizaciones dice usar IaC, pero sólo el 6\pct{} tiene
  todo su parque codificado (P) \autocite{firefly2025}.
& Ambicioso pero razonable: la meta no es la cobertura total. \\
\addlinespace[2pt]
Adopción del portal: 26,5\pct{} de no llegar a 0,60
& Fuera de Spotify, la adopción interna de Backstage se cita en torno al
  10\pct{}, una cifra de origen poco claro (P); en 2025, cerca de la mitad de
  las organizaciones tenía un portal interno (P) \autocite{port2025,backstage}.
& \textbf{Probablemente optimista.} Con una media de 0,60, la probabilidad de
  no pasar G3 sube al 48\pct{} y R7 resta 14,4; el portafolio sigue en 717,6. \\
\addlinespace[2pt]
Primero la plataforma, después el autoservicio
& El 90\pct{} de las organizaciones ya adoptó la ingeniería de plataformas
  \autocite{dora2025}; Gartner preveía que el 80\pct{} tendría equipos de
  plataforma en 2026 \autocite{gartner2023pe}.
& Respalda el orden R3 antes que R7. Es una previsión, no una medición. \\
\addlinespace[2pt]
Sobrecosto medio cercano al 10\pct
& En 1.471 proyectos de tecnología, el sobrecosto medio fue del 27\pct, y uno
  de cada seis tuvo un sobrecosto medio del 200\pct{}
  \autocite{flyvbjerg2011}.
& \textbf{Optimista.} Con un 27\pct, la probabilidad de pasarse sube al
  60\pct{} y el P80 a 630,9. El VAN esperado sigue en 661,8. \\
\bottomrule
\end{tabularx}
\end{table}

El contraste me corrige en dos puntos, y los dos van en la misma dirección: la
adopción de los portales y los sobrecostos son peores en el sector de lo que yo
supuse. En los dos casos el portafolio aguanta gracias a las puertas, pero el
segundo pide una decisión: \textbf{si la dirección quiere un P80 con la tasa de
Flyvbjerg, la reserva tiene que subir de 73,8 a 157,3 miles de dólares}, el
33\pct{} de la base. La otra opción es aceptar una probabilidad de pasarse
cercana al 60\pct. Lo dejo planteado con el número calculado.

\subsection*{Preguntas para un profesional}

\begin{enumerate}
\item ¿Qué fracción de los recursos creados a mano en un parque real no se
  puede importar automáticamente? Es el criterio de G1 y el riesgo K1.
\item ¿Qué adopción voluntaria ha visto en un portal interno seis meses después
  de abrirlo? Es la $\sigma$ de G3.
\item ¿Cuánto cuesta de verdad una auditoría o una excepción manual de
  política? Es lo que falta para saber si R4 se paga solo.
\item ¿Cuál fue el sobrecosto real de sus últimos tres proyectos de plataforma?
  Calibraría la reserva mejor que la media de Flyvbjerg.
\end{enumerate}

\section{Por qué vale la pena}
\label{sec:porque}

Creo que el portafolio vale la pena por tres razones.

\textbf{Se paga, y con margen.} Con 547,4 miles de dólares comprometidos, el VAN
esperado a 2031 es de 727,3, y en el 95\pct{} de los escenarios simulados
supera los 456,4. El flujo acumulado toca fondo a finales de 2027, en $-205{,}1$,
y se vuelve positivo durante 2029. Aun con el sobrecosto medio del sector, el
VAN esperado es de 661,8.

\textbf{Mueve las figuras que la estrategia prometió.} Cuatro de los cinco
compromisos de la hoja de ruta dependen de estos cuatro proyectos: cobertura,
reconstrucción, tiempo de ciclo y cumplimiento. El quinto, la deriva, depende
además de R8.

\textbf{El riesgo está acotado.} Cada pregunta abierta tiene su prototipo y su
puerta antes de que se gaste la mayor parte del dinero. El peor caso de R7, que
es el proyecto más incierto, cuesta 26,5. Es exactamente lo que el libro quiere
decir cuando pide que un proyecto de I+D mida antes de comprometer, y lo que en
mi experiencia con pilotos en la empresa separa los que se pagan de los que se
quedan a medio camino.

\section{Limitaciones}
\label{sec:limitaciones}

\begin{enumerate}
\item \textbf{Las probabilidades son mías.} Las de las puertas (25, 20 y
  26,5\pct) y los márgenes de sobrecosto no salen de una base de datos. El
  contraste con el sector sugiere que las dos más importantes son optimistas.
\item \textbf{El modelo sólo cuenta tres beneficios.} El plazo de salida al
  mercado y el costo de las auditorías no tienen precio en el modelo, y por eso
  R4 y R7 aparecen peor de lo que probablemente son.
\item \textbf{El horizonte de 2031 castiga a R7.} Llega en 2029 y tiene sólo dos
  años completos para devolver su costo. Con 2033 aporta valor positivo.
\item \textbf{R1 se valora desde el 55\pct{} de cobertura de 2026}, que es el
  dato de la hoja de ruta, aunque una parte de esa mejora ya se logró con el
  dinero ya gastado. Eso infla el aporte de R1, pero no cambia ninguna decisión:
  R1 se mantiene en todos los casos.
\item \textbf{No hubo entrevistas.} El contraste es documental, y las cuatro
  preguntas de la sección~\ref{sec:contraste} quedan pendientes.
\end{enumerate}

\printbibliography[title={Referencias}]

\end{document}
